\documentclass[12pt,a4paper]{article}
\usepackage{cmap} % Makes the PDF copiable. See http://tex.stackexchange.com/a/64198/25761
\usepackage[T1]{fontenc}
\usepackage[brazil]{babel}
\usepackage[utf8]{inputenc}
\usepackage{amsmath}
\usepackage{amsfonts}
\usepackage{amssymb}
\usepackage{amsthm}
\usepackage{textcomp} % \degree
\usepackage{gensymb} % \degree
\usepackage[usenames,svgnames,dvipsnames]{xcolor}
\usepackage{hyperref}
\usepackage{graphicx}
\usepackage[margin=2cm]{geometry}

\hypersetup{
    colorlinks = true,
    allcolors = {blue}
}

\newcommand{\fixme}{{\color{red}(...)}}

\newcommand*\sen[1]{\operatorname{sen}{#1}}
\newcommand*\abs[1]{\left|#1\right|}

\newcommand*\tipo{PROVA I}
\newcommand*\turma{NEX152-D}
\newcommand*\disciplina{CDI0001}
\newcommand*\eu{Helder G. G. de Lima}
\newcommand*\data{16/09/2015}

\author{\eu}
\title{\tipo - \disciplina}
\date{\data}

\begin{document}
\thispagestyle{empty}
\newgeometry{margin=2cm,bottom=0.5cm}
\begin{center}
\includegraphics[width=9.0cm]{marca}
\noindent\begin{tabular}{l c c r}
  \textbf{\disciplina}
& \textbf{\tipo}
& \textbf{\data}
& \textbf{\turma}
\end{tabular}
\\ Prof. \eu\footnote{
Este é um material de acesso livre distribuído sob os termos da licença \href{https://creativecommons.org/licenses/by-sa/4.0/deed.pt_BR}{Creative Commons BY-SA 4.0}.}
\end{center}

\noindent Nome do(a) aluno(a): \rule{13cm}{0.01cm}

%\section*{Instruções}
\begin{center}\fbox{
\begin{minipage}{14cm}

{\footnotesize
\begin{itemize}
\renewcommand{\theenumi}{\Roman{enumi}}
\item Identifique-se em todas as folhas.
\item Mantenha o celular e os demais equipamentos eletrônicos desligados durante a prova.
\item Justifique cada resposta com cálculos ou argumentos baseados na teoria estudada.
\item Escolha \textsc{\textbf{uma}} das 6 questões para \textsc{\textbf{não}} fazer (ela não será corrigida): \framebox(30,10){}

\end{itemize}
}

\end{minipage}
}
\end{center}


%\section*{Questões}

\begin{enumerate}
\item Seja $g: D \to \mathbb{R}$ a função dada por $g(x) = \frac{x}{\sqrt{x^2}} + \sqrt{8 - x} - \sqrt{8 + x}$. Explique:
\begin{enumerate}
\item (0,5) Qual é o domínio $D$ em que $g$ está definida?
\item (0,5) A função $g$ é par, ímpar ou nenhuma das duas coisas?
\item (0,5) A função $g$ é contínua nos pontos em que está definida? Por quê?
\item (0,5) Quanto é $\lim\limits_{x\to 8^-} g(x)$?
\end{enumerate}

\item Seja $f(x) = \dfrac{x^2 - 6x - 7}{2x - 14}$ e considere $\varepsilon > 0$.
\begin{enumerate}
\item (1,5) Obtenha $\delta > 0$ (em função de $\varepsilon$) tal que ``se $0 < |x - 7| < \delta$ então $|f(x) - 4| < \varepsilon$''.
\item (0,5) Interprete o significado disso em termos de limites.
\end{enumerate}

\item (2,0) Determine um intervalo $[n, n+1]$ contendo (pelo menos) uma solução positiva da equação $x^4 - 50x = 30$. O que garante que há realmente uma solução no intervalo encontrado?

\item (2,0) Determine se para algum valor de $k$ a função a seguir é contínua em $x = 3$:
\[f(x)=\left\{\begin{array}{ll}
\dfrac{3-x}{x^2-x-6}, & \text{se}\quad x<3\\
-2 k, & \text{se}\quad x = 3 \\
\dfrac{\sen(2x - 6)}{\sen(30 - 10x)}, & \text{se}\quad x > 3.
\end{array}\right.
\]

\item Determine, se existirem:
\begin{enumerate}
\item (1,0) $\lim\limits_{x\to a} \dfrac{\cos(x)-\cos(a)}{x-a}$
\item (1,0) $\lim\limits_{x\to +\infty} \dfrac{\cos^2(2x)}{4-2x}$
\end{enumerate}

\item Calcule, se existirem:
\begin{enumerate}
\item (1,0) $\lim\limits_{x\to 0} \dfrac{e^{9x}-1}{1-e^{3x}} $
\item (1,0) $\lim\limits_{x\to +\infty} \dfrac{\sqrt{16x^2-5}}{4x+4}$
\end{enumerate}
\end{enumerate}

\newpage
\restoregeometry
\section*{Respostas e observações}

\begin{enumerate}
\item \textit{ Seja $g: D \to \mathbb{R}$ a função dada por $g(x) = \frac{x}{\sqrt{x^2}} + \sqrt{8 - x} - \sqrt{8 + x}$. Explique: }
\begin{enumerate}
\item \textit{ Qual é o domínio $D$ em que $g$ está definida? }
\begin{align*}
D
& = \operatorname{Dom}(g) \\
& = \left\{ x \in \mathbb{R} | x^2 \neq 0 \text{ e } 8-x \geq 0 \text{ e } 8+x \geq 0 \right\} \\
& = \left\{ x \in \mathbb{R} | x \neq 0 \text{ e } 8 \geq x \text{ e } x \geq -8  \right\} \\
& = [-8, 8] \setminus\{ 0 \} \\
& = [-8, 0) \cup (0, 8]
\end{align*}

\item \textit{ A função $g$ é par, ímpar ou nenhuma das duas coisas? }

Tem-se
\begin{align*}
g(-x)
& = \frac{-x}{\sqrt{(-x)^2}} + \sqrt{8 - (-x)} - \sqrt{8 + (-x)} \\
& = \frac{-x}{\sqrt{x^2}} + \sqrt{8 + x} - \sqrt{8 - x} \\
& = -\left( \frac{x}{\sqrt{x^2}} - \sqrt{8 + x} + \sqrt{8 - x} \right) \\
& = -g(x).
\end{align*}
Isso significa que $g$ é uma função ímpar.

\item \textit{ A função $g$ é contínua nos pontos em que está definida? Por quê? }

Uma vez que raízes quadradas e polinômios são funções contínuas, e que a soma, subtração, divisão e composição de funções contínuas resulta em funções contínuas em seus domínios, pode-se concluir que a função $g$ também é contínua (pois é uma soma/composição de funções contínuas).

\item \textit{ Quanto é $\lim\limits_{x\to 8^-} g(x)$? }

Levando em conta que $8 \in \operatorname{Dom}(g)$ e que $g$ é contínua, tem-se
\[
\lim\limits_{x\to 8^-} g(x)
= g(8)
= \frac{8}{\sqrt{8^2}} + \sqrt{8 - 8} - \sqrt{8 + 8}
= 1 + 0 - 4 = -3.
\]

\end{enumerate}

\item \textit{ Seja $f(x) = \dfrac{x^2 - 6x - 7}{2x - 14}$ e considere $\varepsilon > 0$. }
\begin{enumerate}
\item \textit{ Obtenha $\delta > 0$ (em função de $\varepsilon$) tal que ``se $0 < |x - 7| < \delta$ então $|f(x) - 4| < \varepsilon$''. }


Observe que
\begin{align*}
\abs{f(x) - 4} < \varepsilon
& \Leftrightarrow \abs{ \frac{x^2 - 6x - 7}{2x - 14} - 4} < \varepsilon
  \Leftrightarrow \abs{ \frac{x^2 - 6x - 7 - 8x + 56}{2x - 14} } < \varepsilon \\
& \Leftrightarrow \abs{\frac{x^2 - 14x + 49}{2(x - 7)}} < \varepsilon
  \Leftrightarrow \abs{\frac{(x - 7)^2}{x - 7}} < 2\varepsilon
  \Leftrightarrow \abs{x - 7} < 2 \varepsilon.
\end{align*}
Portanto, para um $\varepsilon>0$ dado, basta escolher qualquer $\delta \leq 2\varepsilon$ e será verdade que ``se $0 < \abs{x - 7} < \delta$ então $\abs{f(x) - 4} < \varepsilon$''.

\item \textit{ Interprete o significado disso em termos de limites. }

O significado da propriedade anterior é que a função $f(x)$ se aproxima de 4 quando $x$ tende a $7$, isto é, $\lim\limits_{x\to7} f(x) = 4$.

\end{enumerate}

\item \textit{ Determine um intervalo $[n, n+1]$ contendo (pelo menos) uma solução positiva da equação $x^4 - 50x = 30$. O que garante que há realmente uma solução no intervalo encontrado? }

Seja $f(x) = x^4 - 50x$. Atribuindo alguns valores à variável $x$ (positivos, já que o enunciado não pede soluções negativas), observa-se que
\begin{align*}
f(1) & = 1^4 - 50 \cdot 1 = 1 - 50 = -49 < 30 \\
f(2) & = 2^4 - 50 \cdot 2 = 16 - 100 = -84 < 30 \\
\bf{ f(3) } & = 3^4 - 50 \cdot 3 = 81 - 150 = -69 \bf{ < 30 }\\
\bf{ f(4) } & = 4^4 - 50 \cdot 4 = 256 - 200 = 56 \bf{ > 30 }
\end{align*}
Como $f$ é uma função contínua (é um polinômio) e $f(3) < 30 < f(4)$, o teorema do valor intermediário garante que existe algum $c \in (3,4)$ tal que $f(c) = c^4 - 50 \cdot c = 30$, isto é, $c$ é uma solução da equação indicada.

\item \textit{ Determine se para algum valor de $k$ a função a seguir é contínua em $x = 3$:}
\[f(x)=\left\{\begin{array}{ll}
\dfrac{3-x}{x^2-x-6}, & \text{se}\quad x<3\\
-2 k, & \text{se}\quad x = 3 \\
\dfrac{\sen(2x - 6)}{\sen(30 - 10x)}, & \text{se}\quad x > 3.
\end{array}\right.
\]


Para $x<3$ temos o seguinte limite:
\begin{align*}
\lim\limits_{x\to 3^-} f(x)
& = \lim\limits_{x\to 3^-} \frac{3-x}{x^2-x-6}
  = \lim\limits_{x\to 3^-} \frac{3-x}{(x-3)(x+2)}
  = \lim\limits_{x\to 3^-} \frac{-1}{x+2}
  = \frac{-1}{3 +2}
  = \frac{-1}{5}.
\end{align*}

Para $x>3$, temos
\begin{align*}
\lim\limits_{x\to 3^+} f(x)
& = \lim\limits_{x\to 3^+} \frac{\sen(2x - 6)}{\sen(30 - 10x)}
  = \lim\limits_{x\to 3^+} \frac{\sen(2(x - 3)}{\sen(-10(x - 3))} \\
& = \lim\limits_{u\to 0^+} \frac{\sen(2u)}{\sen(-10u)}
  = \lim\limits_{u\to 0^+} \frac{ 2u \frac{ \sen(2u) }{2u} }{ -10u \frac{ \sen(-10u) }{-10u} }
  = \frac{-1}{5} \frac{ \lim\limits_{u\to 0^+} \frac{ \sen(2u) }{2u} }{ \lim\limits_{u\to 0^+} \frac{ \sen(-10u) }{-10u} } \\
& = -\frac{1}{5} \frac{ \lim\limits_{v\to 0^+} \frac{ \sen(v) }{v} }{\lim\limits_{w\to 0^-}  \frac{ \sen(w) }{w} }
  = -\frac{1}{5}.
\end{align*}

Consequentemente, $f$ é contínua em $x=3$ se, e somente se, $-2k = -1/5$, já que neste caso $\lim\limits_{x\to 3} f(x) = -1/5 = k = f(3)$. Portanto, $k = \frac{-1}{-2 \cdot 5} = \frac{1}{10}$.


\item \textit{ Determine, se existirem: }
\begin{enumerate}
\item $\lim\limits_{x\to a} \dfrac{\cos(x)-\cos(a)}{x-a}$
\begin{align*}
\lim\limits_{x\to a} \frac{\cos x-\cos a}{x-a}
& = \lim\limits_{u \to 0} \frac{\cos(a+u)-\cos a}{u}\\
& = \lim\limits_{u \to 0} \frac{\cos a \cos u -\sen a\sen u-\cos a}{u}\\
& = \lim\limits_{u \to 0} \cos a\frac{(\cos u - 1)}{u} + \sen a\frac{\sen u}{u}\\
& = (\cos a) \cdot 0 + (\sen a) \cdot 1
  = \sen a.\\
\end{align*}

\item $\lim\limits_{x\to +\infty} \dfrac{\cos^2(2x)}{4-2x} = \lim\limits_{x\to +\infty} \dfrac{1}{4-2x} \cdot \cos^2(2x) = 0$, pois $\lim\limits_{x\to +\infty} \dfrac{1}{4-2x} = 0$ e além disso $0 \leq \cos^2(2x) \leq 1$, isto é, a função limitada $\cos^2(2x)$ está sendo multiplicada por uma que tende a zero, o que implica que o produto também tende a zero.

\end{enumerate}

\item \textit{ Calcule, se existirem: }
\begin{enumerate}
\item
\begin{align*}
\lim\limits_{x\to 0} \frac{e^{9x}-1}{1-e^{3x}}
& = \lim\limits_{x\to 0} \frac{9x \cdot \frac{e^{9x}-1}{9x}}{3x \cdot \frac{1-e^{3x}}{3x}}
  = 3\lim\limits_{x\to 0} \frac{\frac{e^{9x}-1}{9x}}{(-1)\frac{e^{3x}-1}{3x}}
  = (-3)\frac{\lim\limits_{x\to 0} \frac{e^{9x}-1}{9x}}{\lim\limits_{x\to 0} \frac{e^{3x}-1}{3x}} \\
& = (-3)\frac{\lim\limits_{u\to 0} \frac{e^{u}-1}{u}}{\lim\limits_{v\to 0} \frac{e^{v}-1}{v}}
  = (-3) \cdot \frac{1}{1}
  = -3.
\end{align*}

\item
\begin{align*}
\lim\limits_{x\to +\infty} \frac{\sqrt{16x^2-5}}{4x+4}
& = \lim\limits_{x\to +\infty} \frac{\sqrt{x^2(16-5/x^2)}}{x(4+4/x)}
  = \lim\limits_{x\to +\infty} \frac{\abs{x} \sqrt{16-5/x^2}}{x(4+4/x)} \\
& = \lim\limits_{x\to +\infty} \frac{ \sqrt{16-5/x^2}}{4+4/x}
  = \frac{ \sqrt{16} }{4}
  = \frac{ 4 }{4}
  = 1.
\end{align*}
\end{enumerate}
\end{enumerate}
\end{document}
